%%%%%%%%%%%%%%%%%%%%%%%%%%%%%%%%%%%%%%%%%%%%%%%%%%%%%%%%%%%%%%
%%%%%%%%%%%% MA-0503 Estadística I %%%%%%%%%%%%%%%%%%
%%%%%%%%%%%%%%%%%%%%%%%%%%%%%%%%%%%%%%%%%%%%%%%%%%%%%%%%%%%%%%
\newpage
\subsection*{CA-0303 Estadística I}
\addcontentsline{toc}{subsection}{CA-0303 Estadística I}

\hypertarget{MA0503}{}
\begin{refsection}
\begin{table}[H]
\centering{
\begin{tabular}{|ll|}
\hline
\textbf{Año:} III  & \textbf{Requisitos:} CA-0203 Herramientas de Ciencia de \\ & Datos y  CA-0721 Introducción a la Probabilidad \\ 
\textbf{Ciclo:} I  & \textbf{Correquisitos:} Ninguno \\ 
\textbf{Tipo de curso:} Teórico-práctico & \textbf{Horas semanales:} 5 \\ 
\textbf{Créditos:} 4 & \textbf{Virtualidad:} P, PV \\ \hline
\end{tabular}
}
\end{table}

\subsubsection*{Descripción del curso}

Este es el primer curso del área de estadísticas de la carrera y tiene como requisito que el estudiante conozca sobre la teoría básica de probabilidades por lo que se ubica en el tercer año y primer ciclo. El curso pretende introducir al estudiante a los resultados teóricos y conceptuales principales de la estadística frecuentista y bayesiana para poder realizar su aplicación en problemas de estimación utilizando datos y supuestos paramétricos. Este curso sirve para que el estudiante tenga una base en estadística que le permita entender modelos estadísticos más elaborados en los siguientes cursos.

\subsubsection*{Objetivo general}

Proporcionar al estudiante las principales herramientas teóricas y conceptuales para la aplicación en problemas de inferencia y estimación de parámetros y elaboración de pruebas de hipótesis usando los enfoques frecuentista y bayesiano.

\subsubsection*{Objetivos específicos}

Al finalizar el curso el/la estudiante estará en la capacidad de:

\begin{enumerate}
\item Aplicar técnicas básicas de visualización y síntesis de información. (SC11, SH01.)
\item Comprender las propiedades teóricas a nivel de suficiencia, suficiencia mínima y completitud de un estimador puntual. (SC08.)
\item Estimar puntualmente y por intervalo parámetros en conjuntos de datos bajo modelos de la familia exponencial. (SC08, SH05.)
\item Inferir parámetros de interés en conjunto de datos desde el punto de vista bayesiano bajo modelos de la familia exponencial. (SC08, SH05.)
\item Elaborar pruebas de hipótesis usando los enfoques frecuentista y bayesiano. (SC08, SH05.)
\end{enumerate}

\subsubsection*{Contenidos}

\begin{enumerate}
\item \textbf{Introducción a ciencia de datos}
\begin{enumerate}
\item Ciclo de exploración de datos: transformar, visualizar y modelar.
\item Visualización de datos.
\item Estadísticos de resumen básicos.
\end{enumerate}
\item \textbf{Estadística inferencial frecuentista}
\begin{enumerate}
\item Estimación puntual.
\item Estimación por intervalo.
\end{enumerate}
\item \textbf{Estadística bayesiana}
\begin{enumerate}
\item Familias conjugadas.
\item Intervalos de credibilidad.
\end{enumerate}
\item \textbf{Pruebas de hipótesis}
\begin{enumerate}
\item Pruebas clásicas.
\item Pruebas no paramétricas.
\item Factores de Bayes.
\end{enumerate}
\end{enumerate}

\subsubsection*{Metodología}

Se espera que la persona docente aplique al menos tres de las siguientes estrategias didácticas:

\begin{enumerate}
\item Proyecto.
\item Presentaciones orales a público fuera del grupo.
\item Avances de proyecto usando un caso de estudio.
\item Búsqueda e interpretación de información demográfica en fuentes externas.
\item UVE heurística de Gowin.
\item Foros.
\item Clases magistrales: sincrónicas y asincrónicas.
\item Revisión por pares.
\item Laboratorios guiados en R y en hojas de cálculo.
\end{enumerate}

Se recomienda que la persona docente incluya, a lo largo del desarrollo del curso, ejemplos de variables de frecuencia y severidad, en particular ejemplos de modelación de variables tipo conteo y variables continuas positivas. Además se espera que el/la docente motive la investigación desde la primera semana de clases a través de un proyecto que abarque las 15 semanas de clase en temas relacionados con el ámbito de medición de riesgos e incertidumbre en un rango amplio de aplicación.

\subsubsection*{Evaluación}

\begin{enumerate}
\item Exámenes.
\item Presentaciones orales.
\item Rúbricas.
\item Ensayos cortos de aspectos relacionados con el curso.
\item Proyecto.
\end{enumerate}

\subsubsection*{Bibliografía}

\begin{enumerate}
\item Bickel, P. J. y K. A. Doksum (2000). \emph{Mathematical Statistics: Basic Ideas and Selected Topics}, Vol.\ I (2nd ed.). Prentice Hall.
\item Casella, G. y R. L. Berger (2001). \emph{Statistical Inference.} Cengage Learning.
\item DeGroot, M. H. y M. J. Schervish (2012). \emph{Probability and Statistics} (4th ed.). Pearson Education, Inc.
\item Evans, M. J. y J. S. Rosenthal (2010). \emph{Probability and Statistics. The Science of Uncertainty} (2nd ed.). W.\ H.\ Freeman and Company.
\item Hogg, R. V., J. W. McKean y A. T. Craig (2013). \emph{Introduction to Mathematical Statistics} (7th ed.). Pearson.
\item Hogg, R. V., E. A. Tanis y D. L. Zimmerman (2015). \emph{Probability and Statistical Inference.} Pearson.
\item Rice, J. A. (2007). \emph{Mathematical Statistics and Data Analysis} (3rd ed.). Thomson Brooks-Cole.
\item Robert, C. P. (2007). \emph{The Bayesian Choice: From Decision-Theoretic Foundations to Computational Implementation} (2nd ed.). Springer Texts in Statistics. Springer Verlag, New York.
\item Wasserman, L. (2003). \emph{All of Statistics: A Concise Course in Statistical Inference.} Springer Texts in Statistics. Springer.
\item Wickham, H. y G. Grolemund (2017). \emph{R for data science: import, tidy, transform, visualize, and model data.} O'Reilly Media.
\end{enumerate}

\end{refsection}
